\documentclass[11pt,a4paper]{article}
\usepackage[margin=25mm]{geometry}
\usepackage{amsmath,amssymb,amsthm,booktabs,hyperref}
\newtheorem{theorem}{Theorem}\newtheorem{lemma}[theorem]{Lemma}
\newcommand{\Aut}{\operatorname{Aut}}
\title{Thin geometries from involution triples}
\author{Human Compression project: working application}\date{}
\begin{document}\maketitle
We recover the Suzuki calculation in Section~5 of Jacobs and Leemans,
\emph{An algorithmic analysis of the intersection property}, LMS Journal
of Computation and Mathematics 7 (2004), 284--299. Its input consists of
147 thin, residually connected, residually weakly primitive geometries for
$G=\operatorname{Sz}(8)$. We obtain representatives without importing that
list, and then apply the common shadow criterion. The output is 27 geometries
with the intersection property, 48 with the weak property only, and 72
with neither. These are previously published results.

Throughout, the group preserves types and acts transitively on flags of
each type. Incidence is symmetric and reflexive, with no incident distinct
objects of one type. Every flag extends to a chamber. A geometry is thin
if every flag missing one type has exactly two extensions. Residual
connectedness requires the incidence graph of every residue of rank at
least two to be connected. Residual weak primitivity means that in every
nonempty residue at least one of its actions on a type is primitive.

\begin{lemma}\label{thin}
Let a nonabelian finite simple group $G$ act as above on a thin, residually
connected geometry of rank three. For a chamber $C=\{p_0,p_1,p_2\}$,
write $H_i=G_{p_i}$ and $B=\bigcap_i H_i$. There are distinct involutions
$t_0,t_1,t_2$ such that
\[
 B=1,\qquad G=\langle t_0,t_1,t_2\rangle,\qquad
 H_i=\langle t_j:j\ne i\rangle,
 \qquad H_i\cap H_j=\langle t_k\rangle\quad(\{i,j,k\}=\{0,1,2\}).
\]
\end{lemma}
\begin{proof}
Put $R_i=\bigcap_{j\ne i}H_j$. Flag transitivity on the two extensions
of $C\setminus\{p_i\}$ gives $[R_i:B]=2$, so $B$ is normal in each $R_i$.

The chamber graph is connected. To see this, a rank-two residue has a
connected incidence graph, whose edges are its chambers; successive edges
sharing a vertex give a gallery. In rank three take a path of incident
objects between a vertex of one chamber and a vertex of another. Extend
each edge to a chamber. Consecutive chosen chambers contain a common
object, and the rank-two residue of that object supplies a gallery between
them. The end chambers connect in the same way. Thus all chambers are
connected by changes of one object.

The chambers adjacent to $C$ across its $i$-panel are images of $C$ under
$R_i$. Transporting along a gallery proves
$G=\langle R_0,R_1,R_2\rangle$: the generated subgroup reaches every
chamber, and it contains their common base stabilizer $B$. Therefore
$B\lhd G$. Since $B<R_i$, simplicity gives $B=1$. Each $R_i$ now has a
unique nonidentity element $t_i$, an involution. Within the residue of
$p_i$, the same gallery argument gives
$H_i=\langle R_j:j\ne i\rangle$. Distinctness follows because an element
in two $R_i$ lies in $B$. The stated intersections follow from the
stabilizer definitions.
\end{proof}

\begin{lemma}\label{recognize}
Conversely take three distinct involutions generating $G$, and put
$H_i=\langle t_j:j\ne i\rangle$. Suppose the three intersections in
Lemma~\ref{thin} hold. Their coset incidence system is a flag-transitive
geometry if and only if, for each ordering $i,j,k$ of the types,
\begin{equation}\label{products}
        (H_iH_j)\cap(H_iH_k)=H_i(H_j\cap H_k).
\end{equation}
It then is thin and residually connected. It is residually weakly primitive
if and only if at least one $H_i$ is maximal in $G$ and all three numbers
$|H_i|/2$ are prime.
\end{lemma}
\begin{proof}
Any incident pair of cosets has a common element, so $G$ is transitive
on incident pairs. Normalize the cosets of types $j,k$ to $H_j,H_k$.
The cosets $H_i g$ incident with both correspond exactly to the elements
of $(H_iH_j)\cap(H_iH_k)$. They extend to chambers in the base chamber's
orbit precisely when they meet $H_j\cap H_k$. This gives
\eqref{products}. Every flag then extends to such a chamber. The right
side is contained in the left side; since the triple intersection is
trivial, its size is $2|H_i|$. Thus equality is checked without constructing
an entire incidence graph.

The base chamber stabilizer is trivial and its panel stabilizers have
order two, so the geometry is thin. Each $H_i$ is generated by the two
panel stabilizers within it; the whole group is generated by all three.
Their coset incidence graphs are connected, proving residual connectedness.
Each rank-two residue is the polygon associated to a dihedral group of
order $2m_i$, where $m_i=|H_i|/2$. Its vertex stabilizers are reflections.
A reflection subgroup is maximal in a dihedral group of order $2m_i$
exactly when $m_i$ is prime: intermediate subgroups correspond to proper
divisors of the rotation order. Rank-one residues have degree two and
are primitive. The remaining condition is primitivity on some type of
the whole geometry, equivalent to maximality of its stabilizer $H_i$.
\end{proof}

\section*{A finite enumeration with an explicit orbit reduction}
The ambient permutation group is the previously validated degree-65 model
used in the rank-two Suzuki calculation. We retain that model and GAP's
exact finite group algorithms. The computation checks order 29120 and
simplicity, one conjugacy class containing all 455 involutions, and an
automorphism group of order 87360. Its action on those involutions has
trivial kernel. It also computes representatives of all maximal subgroups,
whose orders are $448,52,20,14$. These are explicit computational premises;
their implementations have not been formalized in Lean here.

Let $\mathcal A$ be the computed automorphism group acting on the 455
involutions, numbered in a fixed order. Fix the first involution. For its
stabilizer $\mathcal A_1$, take one representative $b$ from each orbit on
the other involutions. For each $b$, take one $c$ from each orbit of
$\mathcal A_{1,b}$ on the remaining involutions. This gives all ordered
triples up to $\mathcal A$, exactly once.

To forget the type order, normalize each of the six permutations of a
triple by the same stabilizer procedure and retain the smallest resulting
pair of indices $(b,c)$. The normalization is independent of transporter
choices: after placing the first entry at 1, any two choices differ by
$\mathcal A_1$; after placing the second entry at $b$, they differ by
$\mathcal A_{1,b}$. Thus the final orbit of the third entry is fixed.
Equality of the minimum pairs is equivalent to equivalence under
$\mathcal A$ and a permutation of the three types.

For each retained triple the program checks generation, the exact
subgroup intersections, all three product equalities~\eqref{products},
and the primitivity conditions of Lemma~\ref{recognize}. Top primitivity
is checked by conjugacy to the computed maximal subgroups, not inferred
from a subgroup order. The resulting counts are
\begin{center}\begin{tabular}{lr}\toprule
Stage & Count\\\midrule
Ordered triples modulo $\mathcal A$ & 1083\\
After forgetting the three type names & 195\\
Generating triples & 186\\
Triples satisfying all marked subgroup intersections & 186\\
Flag-transitive thin geometries & 183\\
Residually weakly primitive geometries & 147\\\bottomrule
\end{tabular}\end{center}
The three rejected generating intersection systems show why a subgroup
intersection test alone does not establish the required geometry.

For these thin connected geometries the type-preserving automorphism
group is exactly $G$. Indeed an automorphism fixing one chamber fixes its
unique neighbors across every panel, and then every chamber by connectedness.
The full type-preserving group is consequently semiregular on chambers,
while $G$ is already transitive. It can have no more elements than $G$.
Thus an isomorphism between two such geometries induces an automorphism
of $G$ and a permutation of types, exactly the equivalence used above.

\section*{Shadow tests and their limits}
For each retained triple, the common shadow theorem expresses every
shadow as an orbit of a chamber stabilizer intersection in $H_i\backslash G$.
It checks all normalized object--flag pairs and every allowable witness
flag. The implementation is the earlier reviewed shadow implementation,
with the ambient order and input triples changed. It caches conjugate
coset actions: if $H'=H^g$, the base point changes from $p$ to $p^g$.
The action is unchanged, and the stabilizer of the new point is checked
to equal $H'$. No shadow or witness is removed by this optimization.

The full eligible-extension test is repeated in all three coset actions
for each geometry, providing 1764 flag-extension identities in addition
to the product tests above. Each of the 441 type tests is recorded.
The resulting numbers are 27 with all three types passing, 48 with some
but not all types passing, and 72 with no type passing. All rank-two
residues are polygons of lengths $2m_i$ with $m_i\in\{2,5,7,13\}$.
For $m_i>2$ there is no quadrilateral: shadows of distinct opposite-type
vertices meet in at most one object, which itself is a witness whenever
present. If $m_i=2$, the residue is the complete bipartite graph $K_{2,2}$;
the two opposite-type shadows are the entire target type, with the empty
flag as witness. Flags already containing the target type have singleton
shadows and use that object as witness; the empty input flag has the
whole type as shadow. This proves the rank-two intersection property in
both cases. Rank-one residues are automatic.
Consequently the rank-three answers are also the full residual answers.

The orbit and product arguments enumerate all \emph{thin} geometries
with the stated hypotheses. They do not prove that every residually
weakly primitive rank-three geometry of $\operatorname{Sz}(8)$ is thin.
That separate classification assertion is retained from the earlier
classification cited by Jacobs and Leemans. We claim no recovery of its
unavailable manuscript. The present result is the reported shadow split
on its 147 thin inputs.

The thin chamber argument and stabilizer orbit enumeration are standard
regular-geometry and permutation-group methods. The shadow normalization
and geometric intersection properties are from Jacobs and Leemans.
Our contribution here is an explicit common implementation and an
independently generated list of inputs. Neither the count nor its
mathematical priority is new. This draft awaits independent review; it
adds no original-manuscript credit, numerical compression factor or
complete Lean formalization.
\end{document}
